\section*{अध्याय \ref{ch:b3:photochemistry} --- प्रकाशरसायन}

\begin{solution}{exo:b3:photochemistry:1}
$E = hc/\lambda = \qty{5.442e-19}{J}$; $0.100/\num{5.442e-19} = \num{1.84e17}$ फ़ोटॉन प्रति सेकंड, अर्थात्
\qty{3.05e-7}{einstein.s^{-1}}।
\end{solution}

\begin{solution}{exo:b3:photochemistry:2}
अवशोषित: $\num{1.00e-7} \times 600 = \qty{6.00e-5}{einstein}$; $\Phi = \num{2.0e-5}/\num{6.0e-5} = 0.33$।
\end{solution}

\begin{solution}{exo:b3:photochemistry:3}
अवशोषित फ़ोटॉनों से कहीं अधिक अणु अभिक्रिया करते हैं: प्रत्येक प्रकाश-अपघटित \ce{Cl2} एक
शृंखला आरंभ करता है, $\ce{Cl + H2 -> HCl + H}$ और $\ce{H + Cl2 -> HCl + Cl}$, जो समापन से पहले
हज़ारों चक्कर चलती है।
\end{solution}

\begin{solution}{exo:b3:photochemistry:4}
$D_0(\ce{O2}) = 2 \times 246.79 = \qty{493.58}{kJ/mol}$; $\lambda = hcN_A/D_0 = 0.119627/493\,580 =
\qty{242.4}{nm}$। $\ce{O3 -> O2 + O}$: $246.79 - 145.35 = \qty{101.44}{kJ/mol}$, $\lambda = \qty{1179}{nm}$।
\end{solution}

\begin{solution}{exo:b3:photochemistry:5}
$[Z]/[E] = 22\,000 \times 0.11/(1500 \times 0.40) = 4.0$: 80\,\% $Z$, 20\,\% $E$।
\end{solution}

\begin{solution}{exo:b3:photochemistry:6}\emergencystretch=3em
हेक्सेन-2-ओन: 1,4-द्विमूलक ऐसीटोन के ईनॉल (अतः ऐसीटोन) और
प्रोपीन में विदलित होता है, या 1,2-डाइमेथिल\-साइक्लोब्यूटेन-1-ऑल में बंद हो जाता है। पेन्टेन-2-ओन में $\gamma$
कार्बन अंतस्थ \ce{CH3} है: $\alpha$--$\beta$ आबंध का विदलन $\ce{CH2=CH2}$ छोड़ता है।
\end{solution}

\begin{solution}{exo:b3:photochemistry:7}
$K = \exp[(E_{\mathrm T}(\mathrm D) - 250)/RT]$, जहाँ $RT = \qty{2.478}{kJ/mol}$: 289 से $\eu^{15.7} = \num{7e6}$;
253 से $\eu^{1.2} = 3.4$; 236 से $\eu^{-5.6} = \num{3.5e-3}$। केवल पहला लगभग हर भेंट पर
स्थानांतरण करता है।
\end{solution}

\begin{solution}{exo:b3:photochemistry:8}
$k_2 = \num{6.0e-34} \times (230/300)^{-2.6} = \qty{1.20e-33}{cm^6.s^{-1}}$; $k_2[\mathrm M][\ce{O2}] = \num{1.20e-33} \times
\num{4.0e17} \times \num{8.4e16} = \qty{40}{s^{-1}}$; $[\ce{O}]/[\ce{O3}] = \num{1.0e-3}/40 = \num{2.5e-5}$।
\end{solution}

\begin{solution}{exo:b3:photochemistry:9}
$k(\ce{Cl + O3}) = \num{2.8e-11}\eu^{-250/230} = \num{9.4e-12}$; $k(\ce{Cl + CH4}) = \num{6.6e-12}\eu^{-1240/230} = \num{3.0e-14}$
\unit{cm^3.s^{-1}}। छद्म-प्रथम-कोटि दरें 54 और \qty{0.012}{s^{-1}}: प्रायिकता $54/54.01 = 0.9998$।
\end{solution}

\begin{solution}{exo:b3:photochemistry:10}
$k = \num{2.07e-12}\eu^{-1400/298.15} = \qty{1.89e-14}{cm^3.s^{-1}}$; $[\ce{O3}] = \num{8.0e-3} \times 1.5/\num{1.89e-14} =
\qty{6.3e11}{cm^{-3}}$। \qty{1}{bar} और \qty{298}{K} पर वायु में $p/k_BT = \qty{2.43e19}{cm^{-3}}$ होते हैं: \qty{26}{ppb}। रात में
\ce{NO2} का \omterm{def:b3:photochemistry:photolysis}{प्रकाश-अपघटन} रुक जाता है और NO ओज़ोन को तब तक नष्ट करता है जब तक दोनों में से एक समाप्त न हो जाए।
\end{solution}

\begin{solution}{exo:b3:photochemistry:11}\emergencystretch=3em
$E = hc/\lambda = \qty{6.347e-19}{J}$; $0.050/\num{6.347e-19} = \num{7.88e16}$ फ़ोटॉन प्रति सेकंड $= \qty{1.31e-7}{einstein.s^{-1}}$।
अवशोषित: $\times(1 - 10^{-0.5}) = \times0.684$, \qty{8.95e-8}{einstein.s^{-1}}; दर $0.25 \times \num{8.95e-8} =
\qty{2.24e-8}{mol.s^{-1}}$, अर्थात् \qty{3.0}{mL} में \qty{7.5e-6}{mol.L^{-1}.s^{-1}}।
\end{solution}

\begin{solution}{exo:b3:photochemistry:12}
$I_0 = \num{3.0e-6}/(1.25 \times 60) = \qty{4.0e-8}{einstein.s^{-1}}$। कम रूपांतरण पर प्रकाशमापी
तब भी सारा प्रकाश अवशोषित करता है, और रंगीन उत्पाद उसे छानते नहीं।
\end{solution}

\begin{solution}{pb:b3:photochemistry:1}\emergencystretch=3em
\textbf{1.} $2 \times 246.79 = \qty{493.58}{kJ/mol}$।
\textbf{2.} $\lambda = 0.119627/493\,580 = \qty{242.4}{nm}$।
\textbf{3.} $246.79 - 145.35 = \qty{101.44}{kJ/mol}$; \qty{1179}{nm}।
\textbf{4.} ओज़ोन का दृश्य अवशोषण दुर्बल है; उसका प्रबल पराबैंगनी अवशोषण
(लगभग 200 से \qty{310}{nm}) उस विकिरण को हटा देता है जो DNA और त्वचा को क्षति पहुँचाता है।
\textbf{5.} \qty{242}{nm} से नीचे का प्रकाश नीचे आते हुए स्वयं \ce{O2} द्वारा अवशोषित हो जाता है: उसका थोड़ा ही भाग
ऊपरी समतापमंडल से नीचे पहुँचता है।
\textbf{6.} $\ce{O2 ->[h\nu] 2 O}$; $\mathrm{O + O_2 + M \to O_3 + M}$; $\ce{O3 ->[h\nu] O2 + O}$; $\ce{O + O3 -> 2 O2}$।
\textbf{7.} $k_2 = \qty{1.20e-33}{cm^6.s^{-1}}$; $k_2[\mathrm M] = \qty{4.79e-16}{cm^3.s^{-1}}$।
\textbf{8.} $k_4 = \num{8.0e-12}\eu^{-2060/230} = \qty{1.03e-15}{cm^3.s^{-1}}$।
\textbf{9.} $\num{1.0e-3}/(\num{4.79e-16} \times \num{8.4e16}) = \num{2.49e-5}$।
\textbf{10.} विषम-ऑक्सीजन संतुलन $2j_1[\ce{O2}] = 2k_4[\ce{O}][\ce{O3}]$, जहाँ $[\ce{O}] = j_3[\ce{O3}]/(k_2[\mathrm M][\ce{O2}])$ (प्रमेय
देखिए)।
\textbf{11.} $[\ce{O3}] = \num{8.4e16}\sqrt{\num{1.0e-11} \times \num{4.79e-16}/(\num{1.0e-3} \times \num{1.03e-15})} =
\qty{5.7e12}{cm^{-3}}$, मिश्रण अनुपात $\num{1.4e-5}$ (\qty{14}{ppm})।
\textbf{12.} $\num{2.49e-5} \times \num{5.7e12} = \qty{1.4e8}{cm^{-3}}$।
\textbf{13.} उत्प्रेरकी चक्र (नाइट्रोजन ऑक्साइड, हाइड्रोजन ऑक्साइड, क्लोरीन) विषम ऑक्सीजन के लिए
ह्रास के अतिरिक्त मार्ग जोड़ते हैं।
\textbf{14.} $\ce{Cl + O3 -> ClO + O2}$, $\ce{ClO + O -> Cl + O2}$; नेट $\ce{O + O3 -> 2 O2}$।
\textbf{15.} $1/(\num{9.44e-12} \times \num{5.72e12}) = \qty{0.019}{s}$।
\textbf{16.} $1/(\num{4.033e-11} \times \num{1.423e8}) = \qty{174}{s}$।
\textbf{17.} $174/0.0185 = \num{9.4e3}$।
\textbf{18.} व्यावहारिक रूप से सारा क्लोरीन ClO है: ह्रास $2 \times \num{5.74e-3} \times \num{1.0e7} = \qty{1.1e5}{cm^{-3}.s^{-1}}$,
जबकि चैपमैन चरण के लिए $2k_4[\ce{O}][\ce{O3}] = \qty{1.7e6}{cm^{-3}.s^{-1}}$: चक्र लगभग
7\,\% जोड़ता है।
\textbf{19.} $\ce{ClO + O}$, धीमा चरण (\qty{174}{s})।
\textbf{20.} $53.8/(53.8 + 0.012) = 0.99978$।
\textbf{21.} $\num{5.74e-3}/(\num{5.74e-3} + \num{1.41e-13} \times \num{1.0e9}) = 0.976$।
\textbf{22.} $p = 0.976$।
\textbf{23.} $0.976/0.024 = 40$।
\textbf{24.} बादलों पर $\ce{HCl + ClONO2 -> Cl2 + HNO3}$ क्लोरीन को मुक्त करता है, और नाइट्रिक अम्ल
कणों में रह जाता है: \ce{NO2} के बिना $p_2$ 1 के निकट पहुँचता है और शृंखला
बहुत लंबी हो जाती है; ClO द्वितय से होकर जाने वाला एक चक्र O परमाणुओं के बिना ओज़ोन नष्ट करता है।
\textbf{25.} \textbf{पकड़े जाने से पहले प्रति क्लोरीन परमाणु लगभग 40 ओज़ोन अणु} (\qty{30}{km}
के इस मॉडल में); प्रत्येक परमाणु अपने भंडारों से अनेक बार मुक्त होता है।
\end{solution}
